\documentclass{amsart}
% For dealing with references we use the comment environment
\usepackage{verbatim}
\newenvironment{reference}{\comment}{\endcomment}
%\newenvironment{reference}{}{}
\newenvironment{slogan}{\comment}{\endcomment}
\newenvironment{history}{\comment}{\endcomment}

% For commutative diagrams we use Xy-pic
\usepackage[all]{xy}

% We use 2cell for 2-commutative diagrams.
\xyoption{2cell}
\UseAllTwocells

% We use multicol for the list of chapters between chapters
\usepackage{multicol}

% This is generally recommended for better output
\usepackage{lmodern}
\usepackage[T1]{fontenc}

% For cross-file-references

% Package for hypertext links:
\usepackage{hyperref}

% For any local file, say "hello.tex" you want to link to please

% Theorem environments.
%
\theoremstyle{plain}
\newtheorem{theorem}[subsection]{Theorem}
\newtheorem{proposition}[subsection]{Proposition}
\newtheorem{lemma}[subsection]{Lemma}

\theoremstyle{definition}
\newtheorem{definition}[subsection]{Definition}
\newtheorem{example}[subsection]{Example}
\newtheorem{exercise}[subsection]{Exercise}
\newtheorem{situation}[subsection]{Situation}

\theoremstyle{remark}
\newtheorem{remark}[subsection]{Remark}
\newtheorem{remarks}[subsection]{Remarks}

\numberwithin{equation}{subsection}

% Macros
%
\def\lim{\mathop{\mathrm{lim}}\nolimits}
\def\colim{\mathop{\mathrm{colim}}\nolimits}
\def\Spec{\mathop{\mathrm{Spec}}}
\def\Hom{\mathop{\mathrm{Hom}}\nolimits}
\def\Ext{\mathop{\mathrm{Ext}}\nolimits}
\def\SheafHom{\mathop{\mathcal{H}\!\mathit{om}}\nolimits}
\def\SheafExt{\mathop{\mathcal{E}\!\mathit{xt}}\nolimits}
\def\Sch{\mathit{Sch}}
\def\Mor{\mathop{\mathrm{Mor}}\nolimits}
\def\Ob{\mathop{\mathrm{Ob}}\nolimits}
\def\Sh{\mathop{\mathit{Sh}}\nolimits}
\def\NL{\mathop{N\!L}\nolimits}
\def\CH{\mathop{\mathrm{CH}}\nolimits}
\def\proetale{{pro\text{-}\acute{e}tale}}
\def\etale{{\acute{e}tale}}
\def\QCoh{\mathit{QCoh}}
\def\Ker{\mathop{\mathrm{Ker}}}
\def\Im{\mathop{\mathrm{Im}}}
\def\Coker{\mathop{\mathrm{Coker}}}
\def\Coim{\mathop{\mathrm{Coim}}}

% Boxtimes
%
\DeclareMathSymbol{\boxtimes}{\mathbin}{AMSa}{"02}

%
% Macros for moduli stacks/spaces
%
\def\QCohstack{\mathcal{QC}\!\mathit{oh}}
\def\Cohstack{\mathcal{C}\!\mathit{oh}}
\def\Spacesstack{\mathcal{S}\!\mathit{paces}}
\def\Quotfunctor{\mathrm{Quot}}
\def\Hilbfunctor{\mathrm{Hilb}}
\def\Curvesstack{\mathcal{C}\!\mathit{urves}}
\def\Polarizedstack{\mathcal{P}\!\mathit{olarized}}
\def\Complexesstack{\mathcal{C}\!\mathit{omplexes}}
% \Pic is the operator that assigns to X its picard group, usage \Pic(X)
% \Picardstack_{X/B} denotes the Picard stack of X over B
% \Picardfunctor_{X/B} denotes the Picard functor of X over B
\def\Pic{\mathop{\mathrm{Pic}}\nolimits}
\def\Picardstack{\mathcal{P}\!\mathit{ic}}
\def\Picardfunctor{\mathrm{Pic}}
\def\Deformationcategory{\mathcal{D}\!\mathit{ef}}

\setlength{\emergencystretch}{3em}
\begin{document}
\title[Stacks Project, Tag 0C1L]{571 --- Finite algebraic modifications over algebraically closed fields supported at finitely many points factor into point gluings and tangent-vector contractions}
\maketitle
\noindent Copyright (C) 2005--2025 Johan de Jong. Stacks Project authors. Source: \href{https://stacks.math.columbia.edu/tag/0C1L}{Tag 0C1L}.

\noindent Editorial difficulty note (not author text). Difficulty H1: Induction must construct intermediate scheme algebras while strictly decreasing the finite-length discrepancy. Reduction to a single support point and analysis of the intermediate residue algebra isolate the two minimal algebra extensions, corresponding to gluing points or collapsing a tangent vector..

\noindent Editorial provenance note (not author text). History: excerpt from revision \texttt{a0444\allowbreak 6e57e\allowbreak c1fbc\allowbreak 252a8\allowbreak 71afc\allowbreak ec775\allowbreak 2fb28\allowbreak 07b14}, curves.tex, lines 3039--3103. Reference presentation was adapted; original-author context and core support blocks were retained or added. The mathematical wording is unchanged. Licensed under GNU FDL 1.2 or later; no invariant sections or cover texts. See ../licenses/COPYING. Context blocks reproduce author definitions and supporting results; they are not additional counted problems.

\begin{example}[Glueing points]
\label{example-glue-points}
Let $k$ be an algebraically closed field. Let $f : X' \to X$
be a morphism of algebraic $k$-schemes. We say $X$ is
obtained by glueing $a$ and $b$ in $X'$ if the following are true:
\begin{enumerate}
\item $a, b \in X'(k)$ are distinct points which map to the same
point $x \in X(k)$,
\item $f$ is finite and
$f^{-1}(X \setminus \{x\}) \to X \setminus \{x\}$ is an isomorphism,
\item there is a short exact sequence
$$
0 \to \mathcal{O}_X \to f_*\mathcal{O}_{X'} \xrightarrow{a - b} x_*k \to 0
$$
where arrow on the right sends a local section $h$ of $f_*\mathcal{O}_{X'}$
to the difference $h(a) - h(b) \in k$.
\end{enumerate}
If this is the case, then there also is a short exact sequence
$$
0 \to \mathcal{O}_X^* \to f_*\mathcal{O}_{X'}^*
\xrightarrow{ab^{-1}} x_*k^* \to 0
$$
where arrow on the right sends a local section $h$ of $f_*\mathcal{O}_{X'}^*$
to the multiplicative difference $h(a)h(b)^{-1} \in k^*$.
\end{example}

\begin{example}[Squishing a tangent vector]
\label{example-squish-tangent-vector}
Let $k$ be an algebraically closed field. Let $f : X' \to X$
be a morphism of algebraic $k$-schemes. We say $X$ is
obtained by squishing the tangent vector $\vartheta$ in $X'$
if the following are true:
\begin{enumerate}
\item $\vartheta : \Spec(k[\epsilon]) \to X'$ is a closed immersion
over $k$ such that $f \circ \vartheta$ factors through a point $x \in X(k)$,
\item $f$ is finite and
$f^{-1}(X \setminus \{x\}) \to X \setminus \{x\}$ is an isomorphism,
\item there is a short exact sequence
$$
0 \to \mathcal{O}_X \to f_*\mathcal{O}_{X'} \xrightarrow{\vartheta} x_*k \to 0
$$
where arrow on the right sends a local section $h$ of $f_*\mathcal{O}_{X'}$
to the coefficient of $\epsilon$ in $\vartheta^\sharp(h) \in k[\epsilon]$.
\end{enumerate}
If this is the case, then there also is a short exact sequence
$$
0 \to \mathcal{O}_X^* \to f_*\mathcal{O}_{X'}^*
\xrightarrow{\vartheta} x_*k \to 0
$$
where arrow on the right sends a local section $h$ of $f_*\mathcal{O}_{X'}^*$
to $\text{d}\log(\vartheta^\sharp(h))$ where
$\text{d}\log : k[\epsilon]^* \to k$
is the homomorphism of abelian groups sending $a + b\epsilon$ to $b/a \in k$.
\end{example}

\begin{lemma}
\label{lemma-factor-almost-isomorphism}
Let $k$ be an algebraically closed field. Let $f : X' \to X$ be a
finite morphism of algebraic $k$-schemes such that
$\mathcal{O}_X \subset f_*\mathcal{O}_{X'}$ and such that $f$ is an
isomorphism away from a finite set of points. Then there is a factorization
$$
X' = X_n \to X_{n - 1} \to \ldots \to X_1 \to X_0 = X
$$
such that each $X_i \to X_{i - 1}$ is either the glueing of
two points or the squishing of a tangent vector
(see Examples \ref{example-glue-points} and
\ref{example-squish-tangent-vector}).
\end{lemma}

\begin{proof}
Let $U \subset X$ be the maximal open set over which $f$ is an isomorphism.
Then $X \setminus U = \{x_1, \ldots, x_n\}$ with $x_i \in X(k)$.
We will consider factorizations $X' \to Y \to X$ of $f$ such that
both morphisms are finite and
$$
\mathcal{O}_X \subset g_*\mathcal{O}_Y \subset f_*\mathcal{O}_{X'}
$$
where $g : Y \to X$ is the given morphism. By assumption
$\mathcal{O}_{X, x} \to (f_*\mathcal{O}_{X'})_x$ is an isomorphism
onless $x = x_i$ for some $i$. Hence the cokernel
$$
f_*\mathcal{O}_{X'}/\mathcal{O}_X = \bigoplus \mathcal{Q}_i
$$
is a direct sum of skyscraper sheaves $\mathcal{Q}_i$ supported at
$x_1, \ldots, x_n$.
Because the displayed quotient is a coherent $\mathcal{O}_X$-module,
we conclude that $\mathcal{Q}_i$ has finite length over
$\mathcal{O}_{X, x_i}$. Hence we can argue
by induction on the sum of these lengths, i.e., the length of
the whole cokernel.

\medskip\noindent
If $n > 1$, then we can define an $\mathcal{O}_X$-subalgebra
$\mathcal{A} \subset f_*\mathcal{O}_{X'}$ by taking the inverse
image of $\mathcal{Q}_1$. This will give a nontrivial factorization
and we win by induction.

\medskip\noindent
Assume $n = 1$. We abbreviate $x = x_1$. Consider the finite
$k$-algebra extension
$$
A = \mathcal{O}_{X, x} \subset (f_*\mathcal{O}_{X'})_x = B
$$
Note that $\mathcal{Q} = \mathcal{Q}_1$ is the skyscraper sheaf
with value $B/A$.
We have a $k$-subalgebra $A \subset A + \mathfrak m_A B \subset B$.
If both inclusions are strict, then we obtain a nontrivial
factorization and we win by induction as above.
If $A + \mathfrak m_A B = B$, then $A = B$ by Nakayama, then
$f$ is an isomorphism and there is nothing to prove.
We conclude that we may assume $B = A + \mathfrak m_A B$.
Set $C = B/\mathfrak m_A B$. If $C$ has more than $2$
$k$-subalgebras, then we obtain a subalgebra between $A$
and $B$ by taking the inverse image in $B$. Thus we may
assume $C$ has exactly $2$ $k$-subalgebras. Thus $C = k \times k$
or $C = k[\epsilon]$ by Lemma \href{https://stacks.math.columbia.edu/tag/0C1I}{lemma no in between over k (Tag 0C1I)}.
In this case $f$ is correspondingly the glueing two points or the
squishing of a tangent vector.
\end{proof}
\end{document}
